\documentclass{article}

\usepackage[preprint]{neurips_2026}

\usepackage[utf8]{inputenc}
\usepackage[T1]{fontenc}
\usepackage{hyperref}
\usepackage{url}
\usepackage{booktabs}
\usepackage{amsfonts}
\usepackage{amsmath}
\usepackage{amssymb}
\usepackage{nicefrac}
\usepackage{microtype}
\usepackage{xcolor}
\usepackage{graphicx}
\usepackage{multirow}
\usepackage{array}
\usepackage{enumitem}

\title{Benchmarking Subset Aggregation for Classical Causal Discovery Under Marginalization Error}

\author{Anonymous Author(s)}

\newcommand{\method}{CompatExp}
\newcommand{\pcstable}{PC-Stable}
\newcommand{\cpdag}{CPDAG}
\newcommand{\forate}{false-orientation rate}

\begin{document}

\maketitle

\begin{abstract}
Subset-estimate-aggregate pipelines are attractive for causal discovery because they recycle small local structure-learning problems into one global graph, but their soundness is fragile when local graphs are learned on observed-variable subsets of a larger system. We study this issue in a fixed-bank benchmark rather than proposing a new discovery algorithm. Across $48$ synthetic datasets spanning Erd\H{o}s--R\'enyi and scale-free graphs, four structural regimes, two sample sizes, and three seeds, we keep one bank of $40$ local \cpdag s per dataset fixed and vary only weighting and merge rules. The benchmark first audits subset distortion against induced-subgraph truth: non-adjacency claims are usually valid ($0.980$--$0.993$ by regime), but direction claims are rarely valid ($0.110$--$0.158$) and are abstained on $0.888$--$0.900$ of opportunities. Under this distortion, compatibility weighting yields only a small hard-regime median reduction in \forate\ over matched uniform weighting ($0.0286$ for \method+$\,$DetThreshold), far below the pre-registered $0.10$ target and statistically null ($p=0.351$). Full-graph baselines remain much stronger: GES reaches mean SHD $6.83$ and mean orientation precision $0.748$, whereas the best subset variants remain near SHD $18.4$--$18.5$ and orientation precision $0.417$--$0.426$. The main conclusion is negative but useful: in this low-compute setting, subset compatibility is not enough to overcome marginalization error, and benchmark design choices such as subset size matter more than the proposed weighting heuristic.
\end{abstract}

\section{Introduction}

Many recent causal discovery papers use internal agreement signals to select, evaluate, or aggregate candidate graphs without access to ground truth \citep{viinikka2018intersection,faller2024selfcompatibility,schkoda2025lovo}. A natural extension is to learn local graphs on overlapping subsets of variables and then merge their symbolic claims into one global \cpdag. That pattern already appears in subset-estimate-aggregate pipelines \citep{wu2025sample}, in bagging-style graph aggregation \citep{wang2014dagbag,debeire2024bootstrap}, and in newer consistency-guided aggregation work \citep{ninad2025vector}. The unresolved issue is not whether graph aggregation can help in general, but whether compatibility-style weighting is actually useful when local graphs come from arbitrary observed-variable subsets of one larger system.

That setting is methodologically dangerous. A \cpdag\ learned on a subset of observed variables need not match the \cpdag\ of the induced subgraph on that subset. Omitted observed variables can create effective confounding, alter collider patterns, and break the transport of local orientations. If so, disagreement among local graphs may reflect structural distortion rather than finite-sample noise. Our paper therefore studies subset aggregation as an empirical benchmark with an explicit soundness audit, not as a new causal discovery principle.

We build on the proposal in this workspace and ask a narrower question: when one fixed bank of local subset \cpdag s is held constant, do compatibility-based weights improve global recovery relative to uniform weighting and a non-falsification bootstrap-stability control? The answer in our benchmark is mostly no. The subset audit shows that local non-adjacency statements are usually correct, but adjacency, direction, and collider claims are heavily distorted. Under these conditions, compatibility weighting produces only a $0.0286$ hard-regime median reduction in \forate\ over matched uniform weighting, with a paired Wilcoxon $p$-value of $0.351$. At the same time, strong full-graph baselines such as GES and \pcstable\ remain far better than every subset variant we tested.

Our contributions are:
\begin{itemize}[leftmargin=*,topsep=2pt,itemsep=1pt]
  \item We define a reproducible fixed-bank benchmark for overlapping-subset \cpdag\ aggregation under explicit marginalization error, covering $48$ synthetic datasets with CPU-only execution and shared local evidence across all subset methods.
  \item We audit subset validity before interpreting aggregation gains. Across regimes, direction validity is only $0.110$--$0.158$ and collider validity only $0.088$--$0.140$, even though non-adjacency validity remains $0.980$--$0.993$.
  \item We show that compatibility weighting is not a robust remedy in this setting: its observed gains over uniform weighting are small, statistically null, and substantially weaker than improvements obtained by changing subset size or using full-graph baselines.
\end{itemize}

Section~\ref{sec:related} situates the benchmark in prior work on internal validation and graph aggregation. Section~\ref{sec:method} defines the fixed-bank aggregation setup and the subset-validity audit. Section~\ref{sec:experiments} presents the benchmark, main results, and ablations. Section~\ref{sec:discussion} discusses implications and limitations, and Section~\ref{sec:conclusion} concludes.

\section{Related Work}
\label{sec:related}

\paragraph{Internal validation for causal discovery.}
Agreement-based evaluation without ground truth predates recent causal discovery benchmarks. Intersection-Validation \citep{viinikka2018intersection} uses overlap agreement to evaluate structure learning, while Self-Compatibility \citep{faller2024selfcompatibility} turns subset disagreement into an internal reliability signal. Leave-One-Variable-Out extends this line toward model selection rather than aggregation \citep{schkoda2025lovo}. Our benchmark borrows the idea that disagreement can be informative, but tests it in a harder setting where local graphs are learned on fixed observed-variable subsets and then merged.

\paragraph{Subset aggregation and consistency-guided merging.}
SEA shows that sample-estimate-aggregate pipelines can be effective at larger scale with learned aggregation modules \citep{wu2025sample}. Ninad et al.\ study consistency-guided aggregation for vector-valued variables and explicitly analyze failure modes when aggregation maps are invalid \citep{ninad2025vector}. Those works substantially narrow novelty for any new subset-aggregation method. Our contribution is therefore empirical and narrower: a controlled benchmark that isolates weighting and merging decisions while auditing the local evidence for marginalization-induced distortion.

\paragraph{Classical graph aggregation and competing full-graph baselines.}
Graph bagging is longstanding in causal structure learning. DAGBag aggregates bootstrap DAGs to reduce finite-sample false positives \citep{wang2014dagbag}, and similar ideas appear in time-series causal discovery \citep{debeire2024bootstrap}. We also compare against standard full-graph learners and selectors: \pcstable\ \citep{colombo2014pcstable}, GES \citep{chickering2002ges}, NOTEARS \citep{zheng2018notears}, and self-compatibility-based model selection \citep{faller2024selfcompatibility}. AutoML-style baseline selection is another alternative to local aggregation \citep{chan2024autocd}, and \citet{malinsky2024cautious} argues for conservatism in constraint-based model selection. Unlike these approaches, our benchmark fixes the local subset evidence and asks only whether compatibility-based weighting adds value once marginalization error is already present.

\section{Method}
\label{sec:method}

\subsection{Fixed-Bank Benchmark Setup}

This paper does not propose a new discovery algorithm. The object of study is a family of symbolic aggregators that all operate on the same precomputed bank of local \cpdag s. For each dataset, the benchmark records $20$ overlapping subsets of size $k=8$ and learns two bootstrap \cpdag s per subset using \pcstable\ with $\alpha=0.01$, yielding $40$ local graphs per dataset. The executed environment used dataset seeds $\{11,23,37\}$, bootstrap seeds $\{101,131\}$, subset-bank seed $509$, two CPU workers, and no GPUs.

All subset methods therefore share identical local evidence. Differences between methods come only from how local graphs are weighted and how weighted claims are merged. This isolates the contribution of compatibility weighting from confounds due to subset construction, local learning, or stochastic variation in the evidence bank.

\subsection{Claim Extraction}

From each local graph $g$ on subset $S_g$, we extract only explicit symbolic claims:
\begin{align*}
  \texttt{adj}(i,j), \quad
  \texttt{nonadj}(i,j), \quad
  \texttt{dir}(i \rightarrow j), \quad
  \texttt{coll}(i \rightarrow k \leftarrow j).
\end{align*}
Undirected edges contribute adjacency evidence only. Missing orientations are treated as abstentions rather than contradictions.

\subsection{Compatibility Weighting}

For two local graphs $g_a$ and $g_b$ with overlap $O = S_a \cap S_b$, we compare only claims whose variables lie in $O$. Contradictions are counted only when both graphs make explicit opposing statements, namely adjacency versus non-adjacency, opposite directions on the same edge, orientation versus non-adjacency, or incompatible collider statements on the same triple. Let $\mathrm{contr}(g_a,g_b)$ denote the number of explicit contradictions and $\mathrm{comp}(g_a,g_b)$ the number of comparable claims. We define the overlap contradiction rate
\[
r(g_a,g_b)=
\begin{cases}
\mathrm{contr}(g_a,g_b)/\mathrm{comp}(g_a,g_b), & \mathrm{comp}(g_a,g_b)>0,\\
0, & \mathrm{comp}(g_a,g_b)=0.
\end{cases}
\]
For each local graph $g$, the compatibility score is the mean contradiction rate against overlapping neighbors,
\[
C_g = \frac{1}{|\mathcal{N}(g)|}\sum_{h \in \mathcal{N}(g)} r(g,h).
\]
Our primary weighting rule is
\[
w_g = \exp(-2C_g).
\]
We compare this \method\ rule against uniform weighting and a bootstrap-stability control that scores each subset using agreement between its two bootstrap graphs. In the main benchmark, the median effective sample size proxy is $40.0$ for uniform weighting and $39.40$ for \method, while median weighted opposition falls from $4.0$ to $3.056$.

\subsection{Merge Rules}

Aggregators combine weighted support tables and then apply deterministic consistency checks to produce a final \cpdag. The main paper studies two merge rules.

\paragraph{Wilson-style screen.}
For a binary claim family, let
\[
\hat{p} = \frac{\text{support}}{\text{support}+\text{opposition}},
\qquad
n_{\mathrm{eff}} = \frac{(\sum_g w_g)^2}{\sum_g w_g^2}.
\]
A claim is accepted if the one-sided $90\%$ Wilson lower bound exceeds $0.5$. Because local graphs overlap heavily, $n_{\mathrm{eff}}$ is only a shrinkage heuristic and not a calibrated sample size.

\paragraph{Deterministic thresholding.}
DetThreshold accepts a claim whenever the weighted support ratio exceeds $0.6$. We also study deterministic rank-based ablations, but our main comparisons use Wilson and DetThreshold because both were pre-registered across all $48$ datasets.

\subsection{Subset-Validity Audit}

The benchmark interprets aggregation only after auditing the local evidence against induced-subgraph truth. For every sampled subset $S$, we compare each learned local graph against the true \cpdag\ of the induced subgraph on $S$. This yields per-claim validity, contradiction, and abstention rates before any aggregation takes place. The audit is essential because compatibility weighting can only be meaningful if local disagreement is at least partially tracking finite-sample error rather than mostly reflecting subset-induced distortion.

\section{Experiments}
\label{sec:experiments}

\subsection{Benchmark and Evaluation Protocol}

The executed benchmark contains $48$ datasets from the Cartesian product of two graph families (Erd\H{o}s--R\'enyi and scale-free), four regimes (linear Gaussian, nonlinear ANM, near-unfaithful linear, and mild misspecification), two sample sizes ($n=200$ and $n=1000$), and three seeds. Every dataset has $p=20$ observed variables. The benchmark manifest confirms that true \cpdag\ edge counts range from $15$ to $23$ in the Erd\H{o}s--R\'enyi family and are $19$ in every scale-free instance.

We evaluate five full-graph baselines (PC, GES, NOTEARS-L1, DAGBag-PC, and SC-Select) and six fixed-bank subset methods (Uniform, BootstrapStability, and \method\ combined with Wilson or DetThreshold merging). Metrics include SHD, skeleton F1, orientation precision, \forate, fraction of recovered adjacencies left undirected, runtime, and peak memory. The environment ran with Python $3.12.3$, two CPU workers, and no GPUs; the timing pilot projected total benchmark cost at $0.206$ hours and did not trigger any fallback reductions.

\subsection{Subset Audit: Local Evidence Is Structurally Distorted}

Table~\ref{tab:audit} summarizes the subset-validity audit. The picture is asymmetric. Non-adjacency claims are highly reliable in every regime, with validity between $0.9798$ and $0.9932$. By contrast, adjacency validity stays near chance at $0.4811$--$0.5156$, direction validity is only $0.1101$--$0.1585$, and collider validity is only $0.0882$--$0.1404$. The local bank is also highly abstaining on orientation-sensitive claims: direction abstention is $0.8881$--$0.9004$ and collider abstention is $0.9853$--$0.9875$.

These numbers sharply constrain what aggregation can plausibly recover. If local direction claims are mostly wrong or absent, compatibility weighting has little trustworthy orientation signal to amplify. Figure~\ref{fig:subset_validity_gain} makes the same point at the dataset level: the correlation between direction validity and the \method-minus-Uniform change in \forate\ is weak (Pearson $-0.066$, Spearman $-0.218$).

\begin{table}[t]
\caption{Subset-validity audit against induced-subgraph truth. Each number is the regime-level validity or abstention rate computed from local graphs before aggregation. High non-adjacency validity contrasts with poor orientation validity, showing that the local bank is structurally distorted in exactly the claim families subset aggregation most needs.}
\label{tab:audit}
\centering
\small
\begin{tabular}{lccccc}
\toprule
Regime & Adj val. & Nonadj val. & Dir val. & Coll. val. & Dir abst. \\
\midrule
Linear Gaussian & 0.5094 & 0.9932 & 0.1537 & 0.1319 & 0.8881 \\
Mild misspecification & 0.5156 & 0.9883 & 0.1473 & 0.1275 & 0.8904 \\
Near-unfaithful linear & 0.4811 & 0.9798 & 0.1585 & 0.1404 & 0.8987 \\
Nonlinear ANM & 0.4906 & 0.9815 & 0.1101 & 0.0882 & 0.9004 \\
\bottomrule
\end{tabular}
\end{table}

\subsection{Main Results: Compatibility Weighting Does Not Rescue Subset Aggregation}

Table~\ref{tab:main} reports the main benchmark summary. Among subset methods, the differences are minimal. Uniform+DetThreshold achieves mean SHD $18.38$, mean skeleton F1 $0.7325$, mean orientation precision $0.4264$, and mean \forate\ $0.5736$. The corresponding \method+DetThreshold values are SHD $18.52$, skeleton F1 $0.7307$, orientation precision $0.4173$, and \forate\ $0.5827$. The Wilson variants are nearly identical to their deterministic counterparts.

The pre-registered hard-regime test leads to the same conclusion. On the near-unfaithful and mild-misspecification slice, \method+DetThreshold reduces the median \forate\ by only $0.0286$ relative to Uniform+DetThreshold. That effect falls well short of the pre-registered $0.10$ success criterion and is statistically null (paired Wilcoxon $p=0.351$, median paired difference $0.0$, bootstrap CI $[0.0, 0.0]$). Against BootstrapStability+DetThreshold, the difference is again $0.0$ with $p=0.638$.

More importantly, all subset variants are substantially worse than strong full-graph baselines. GES is the best overall method with mean SHD $6.83$, mean skeleton F1 $0.9047$, mean orientation precision $0.7477$, and mean \forate\ $0.2523$. PC and DAGBag-PC are close behind. Relative to \method+DetThreshold, the paired median \forate\ difference is $0.2752$ for GES, $0.2638$ for PC, and $0.2344$ for DAGBag-PC, all with $p<10^{-7}$.

\begin{table*}[t]
\caption{Main benchmark summary over $48$ datasets. Best values are in \textbf{bold}. Higher is better for skeleton F1 and orientation precision; lower is better for SHD, false-orientation rate, and runtime. The subset methods are tightly clustered, while full-graph baselines dominate them on every structural metric.}
\label{tab:main}
\centering
\small
\resizebox{\textwidth}{!}{%
\begin{tabular}{lccccc}
\toprule
Method & SHD $\downarrow$ & Skeleton F1 $\uparrow$ & Orient. Prec. $\uparrow$ & False-orient. $\downarrow$ & Runtime (s) $\downarrow$ \\
\midrule
PC & 7.29 & 0.8967 & 0.7239 & 0.2761 & \textbf{0.3409} \\
GES & \textbf{6.83} & \textbf{0.9047} & \textbf{0.7477} & \textbf{0.2523} & 11.9116 \\
NOTEARS-L1 & 18.83 & 0.0000 & 0.0000 & 0.0000 & 0.5132 \\
DAGBag-PC & 8.10 & 0.8839 & 0.7008 & 0.2992 & 2.8400 \\
SC-Select & 8.04 & 0.8840 & 0.6832 & 0.3168 & 1.4722 \\
\midrule
Uniform+Wilson & 18.38 & 0.7325 & 0.4264 & 0.5736 & \textbf{0.0100} \\
Uniform+DetThreshold & 18.38 & 0.7325 & 0.4264 & 0.5736 & 0.0139 \\
BootstrapStability+Wilson & 18.46 & 0.7313 & 0.4179 & 0.5821 & 0.0163 \\
BootstrapStability+DetThreshold & 18.48 & 0.7309 & 0.4175 & 0.5825 & 0.0165 \\
\method+Wilson & 18.50 & 0.7309 & 0.4173 & 0.5827 & 0.0261 \\
\method+DetThreshold & 18.52 & 0.7307 & 0.4173 & 0.5827 & 0.0261 \\
\bottomrule
\end{tabular}
}
\end{table*}

\subsection{Ablations: Subset Size Matters More Than Compatibility Weighting}

Figure~\ref{fig:ablations} and the ablation summaries reinforce the negative main result.

\paragraph{Alternative compatibility weights.}
On the hard-slice ablation, CompatRank+DetThreshold reaches median SHD $17.0$ and median \forate\ $0.6583$, while CompatTopHalf+DetThreshold reaches median SHD $17.5$ and median \forate\ $0.7450$. Neither variant improves on the main \method\ rule in a meaningful way.

\paragraph{Claim-family ablations.}
Removing collider claims is mostly neutral: CompatExp+DetThreshold+NoColliders yields median SHD $16.5$ and median \forate\ $0.6049$. By contrast, removing directional information is destructive. The adjacency-only version has median SHD $20.5$, mean orientation precision $0.1279$, mean \forate\ $0.7888$, and fraction undirected $0.4893$. In other words, the subset methods still need orientation claims even though those claims are noisy.

\paragraph{Contradiction definition.}
Treating one-sided abstention as weak disagreement modestly helps on the hard slice, improving mean orientation precision from $0.4173$ in the main method to $0.4399$ and lowering mean \forate\ from $0.5827$ to $0.5601$. This suggests that the explicit-claims-only contradiction rule is not uniquely responsible for the observed behavior.

\paragraph{Subset-size sensitivity.}
Subset size has the largest effect of any ablation. At $k=10$, Uniform+DetThreshold improves to mean SHD $12.42$, mean skeleton F1 $0.8051$, mean orientation precision $0.5497$, and mean \forate\ $0.4503$; \method+DetThreshold is nearly identical with SHD $12.54$ and \forate\ $0.4524$. At $k=6$, both methods deteriorate sharply: Uniform+DetThreshold has mean SHD $21.63$ and mean \forate\ $0.6302$, while \method+DetThreshold has mean SHD $21.96$ and mean \forate\ $0.6412$. This is a central result of the paper: subset design choices dominate the weighting choice.

\begin{figure}[t]
\centering
\includegraphics[width=\linewidth]{figures/figure1_false_orientation_boxplot.png}
\caption{Distribution of false-orientation rates across regimes and sample sizes. The paired links between Uniform+DetThreshold and \method+DetThreshold show that most dataset-level differences are small.}
\label{fig:false_orientation}
\end{figure}

\begin{figure}[t]
\centering
\includegraphics[width=\linewidth]{figures/figure2_subset_validity_vs_gain.png}
\caption{Subset direction-validity rate versus the dataset-level change in false-orientation rate from replacing Uniform+DetThreshold with \method+DetThreshold. The weak negative correlations (Pearson $-0.066$, Spearman $-0.218$) indicate no strong evidence that compatibility weighting becomes useful exactly when local orientation validity improves.}
\label{fig:subset_validity_gain}
\end{figure}

\begin{figure}[t]
\centering
\includegraphics[width=\linewidth]{figures/figure3_ablations.png}
\caption{Merger and ablation summary. The broad pattern is stable: changing the merge rule or compatibility variant does little, whereas subset size and the availability of directional claims move performance substantially more.}
\label{fig:ablations}
\end{figure}

\begin{figure}[t]
\centering
\includegraphics[width=\linewidth]{figures/figure4_runtime_vs_quality.png}
\caption{Runtime versus quality trade-off. Subset aggregation is very fast at merge time, but its structural quality remains far behind the strongest full-graph baselines. All reported runs used two CPU workers and zero GPUs.}
\label{fig:runtime_quality}
\end{figure}

\section{Discussion and Limitations}
\label{sec:discussion}

The benchmark supports a constrained conclusion. Compatibility weighting is not harmful in the sense of catastrophic failure, but it is also not a reliable solution to the subset-marginalization problem. The local evidence bank already contains severe orientation distortion: direction validity is below $0.16$ in every regime, collider validity below $0.15$, and abstention on those claim families is extremely high. Under those conditions, the weighted merger mostly shuffles weak or corrupted orientation evidence rather than extracting robust global signal.

The strongest practical lesson is that benchmark geometry matters more than the proposed weighting heuristic. Moving from $k=8$ to $k=10$ improves both uniform and compatibility-weighted subset methods by roughly six SHD points and lowers mean \forate\ by about $0.12$--$0.13$, whereas switching from uniform to compatibility weighting at fixed $k=8$ barely changes the outcome. This suggests that future work on subset aggregation should prioritize subset construction, local learner fidelity, and marginalization-aware claim selection over increasingly elaborate weighting rules.

The paper also has several limitations. First, it is a synthetic benchmark only; real data are intentionally out of scope because they would weaken the ground-truth claims central to the soundness audit. Second, the subset methods all use \pcstable\ local learners and symbolic merging, so the conclusions do not automatically transfer to learned aggregators such as SEA. Third, the Wilson screen uses an overlap-sensitive effective sample size heuristic rather than a calibrated uncertainty estimate. Fourth, the planned LOVO baseline was not run because the scoring recipe was not operationally unambiguous without adding unregistered design choices. Finally, the benchmark fixes $p=20$ and a low-compute CPU setting, so larger problems may show different trade-offs.

\section{Conclusion}
\label{sec:conclusion}

We presented a fixed-bank benchmark for subset aggregation in classical causal discovery under explicit marginalization error. The benchmark first audits whether local subset claims are trustworthy relative to induced-subgraph truth and then asks whether compatibility weighting improves global recovery once the local evidence bank is held fixed.

The main findings are negative but informative. Local non-adjacency claims are usually valid ($0.980$--$0.993$), but local direction claims are rarely valid ($0.110$--$0.158$) and highly abstaining. In that regime, compatibility weighting produces only a $0.0286$ hard-regime median reduction in \forate\ over uniform weighting, far below the pre-registered success threshold and statistically null. Meanwhile, full-graph baselines such as GES, PC, and DAGBag-PC remain much stronger, with mean SHD between $6.83$ and $8.10$ versus approximately $18.4$--$18.5$ for the subset methods.

The broader implication is that overlap compatibility is not a substitute for valid local evidence. In low-compute fixed-variable subset aggregation, marginalization error appears to be the dominant bottleneck, and benchmark choices such as subset size matter more than the weighting heuristic itself. Future work should focus on marginalization-aware local representations, better subset design, or aggregation rules that can explicitly reason about omitted-variable distortion instead of treating local \cpdag s as directly transportable evidence.

\bibliographystyle{plainnat}
\bibliography{refs}

\end{document}
